% Generated from locale/tr/content/normal-modal-logic/filtrations/introduction.tex; do not hand-edit.


\olfileid[tr]{nml}{fil}{int}

\olsection{Giriş}

Bir mantıkla ilgili önemli sorulardan biri, onun karar verilebilir olup
olmadığıdır; yani ``bu !!{formula} geçerli mi?'' sorusunu yanıtlayacak etkili
bir yordam bulunup bulunmadığıdır. Önerme mantığı karar verilebilirdir: bir
!!a{formula}{}ün totoloji olup olmadığını doğruluk tablosu kurarak etkili
biçimde sınayabiliriz ve verilmiş bir !!{formula} için doğruluk tablosu
sonludur. Fakat sonsuz sayıda modal model olduğundan, modal bir formülün
bütün modellerde doğru olup olmadığını açık bir biçimde sınayamayız. Verilmiş
  verilmiş bir formülle ilgili bütün sonlu modelleri listeleyebiliriz; çünkü yalnızca
!!{formula} içinde gerçekten geçen !!{propositional variable}s için dünya
alt kümelerinin ataması önemlidir. Erişilebilirlik bağıntısı sabitse olası
farklı $V(p)$ atamaları $W$'nin bütün alt kümeleridir; $\card{W}=n$ ise
bunlardan $2^n$ tane vardır. $!A$ !!{formula}{}ümüz $m$ !!{propositional
variable}s içeriyorsa $n$ dünyalı $2^{nm}$ farklı model vardır. Bunların her
  birinde $!A$'nın bütün dünyalarda doğru olup olmadığını, her dünyadaki
  doğruluk değerini hesaplayarak sınayabiliriz. Elbette bütün olası
erişilebilirlik bağıntılarını da denetlemeliyiz; ancak $n$ dünya üzerinde de
yalnızca sonlu sayıda bağıntı vardır (özellikle $W\times W$'nin alt kümeleri
kadar, yani $2^{n^2}$ tane).

  \Log{K} mantığıyla değil, bir modeller sınıfıyla tanımlanmış bir mantıkla
(örneğin yansımalı geçişli modellerle) ilgileniyorsak, erişilebilirlik
bağıntısının doğru türden olup olmadığını da sınayabilmeliyiz. İlgilendiğimiz
çerçeveler modal !!{formula}s ile tanımlanabildiğinde bunu yapabiliriz
(örneğin \Ax{T} ile \Ax{4}'ün çerçevede geçerli olup olmadığını sınayarak).
O hâlde fikir, bütün sonlu çerçeveleri sırayla ele almak, her birinin
ilgilendiğimiz sınıftan olup olmadığını sınamak, sonra o çerçeve üzerindeki
bütün olası modelleri listeleyip $!A$'nın her birinde doğru olup olmadığını
denetlemektir. Doğru değilse dururuz: $!A$, ilgilenilen modeller sınıfında
geçerli değildir.

Bu fikrin bir sorunu vardır: aramayı ne zaman durdurabileceğimizi bilmiyoruz.
!!{formula}{}ün sonlu bir karşı-modeli varsa yordamımız onu bulur. Fakat sonlu
karşı-modeli yoksa yanıt alamayız. !!{formula} geçerli olabilir (hiç
karşı-modeli yoktur) veya yalnızca sonsuz bir karşı-modeli olabilir; buna
hiç bakmayacağız. Karşı-modeli olan her !!{formula}{}ün sonlu bir karşı-modeli
de olduğunu gösterebilirsek bu sorun aşılır. Bu durumda mantığın
\emph{sonlu model özelliğine} sahip olduğunu söyleriz.

Peki bir mantığın sonlu model özelliğine sahip olduğunu nasıl gösterebiliriz?
Bunun bir yolu, $!A$'nın sonsuz bir (karşı-)modelini sonlu bir modele
dönüştürme yolu bulmaktır. Bu yapılabiliyorsa $!A$'nın doğru olmadığı bir
model bulunduğunda ortaya çıkan sonlu model de $!A$'yı doğru yapmaz. Bu
sonlu model bütün sonlu modeller listemizde görünür ve geçerli olmayan her
formülün geçersiz olduğunu sonunda belirleriz. Formül geçerliyse yordamımız
sona ermez. Ek olarak, yordamımızın verdiği sonlu modelin yalnızca $!A$
!!{formula}{}üne bağlı bir azami büyüklüğü olduğunu gösterebilirsek,
karşı-model ararken sonlu modelleri hangi büyüklüğe kadar sınamamız
gerektiğini biliriz. O noktaya kadar karşı-model bulamadıysak hiç yoktur.
  Yordamımız o zaman her !!{formula} $!A$ için ``$!A$ geçerli mi?'' sorusuna
  gerçekten karar verir.

Sonsuz yapıları sonlu yapılara dönüştürmekte sık işe yarayan bir strateji,
yapının ilgili bakımlardan aynı davranan !!{element}s{}ını
``özdeşleştirmektir.'' $\mModel{M}$ içinde ilgili bakımlardan aynı davranan
sonsuz sayıda dünya varsa, \emph{bütün} dünyaları bu tür dünyaların sonlu
sayıda (kendileri sonsuz olabilen) ``sınıfında'' toplayabiliriz. Başka bir
deyişle dünyalar kümesini doğru biçimde bölümlemeliyiz: her bölüm sonsuz
sayıda dünya içerebilir, ama yalnızca sonlu sayıda bölüm bulunmalıdır. Sonra
dünyaları bu bölümler olan yeni bir $\mModel{M^*}$ modeli tanımlarız. Eski
modelde sonlu sayıda bölüm, yeni modelde sonlu sayıda dünya; yani sonlu bir
model verir. Bir $w$ dünyasının bulunduğu bölüme $[w]$ diyelim. Eski model
$!A$'nın bir karşı-modeliyse yenisinin de öyle olmasını istediğimizden
$\mSat{M}{!A}[w]$ ancak ve ancak $\mSat{M^*}{!A}[{[w]}]$ olmalıdır. Bunun
için hem bölümlemeyi hem $\mModel{M^*}$ modelinin erişilebilirlik bağıntısını
doğru biçimde tanımlamalıyız.

  Bunun nasıl işleyeceğini görmek için önce erişilebilirlik bağıntısının
  tümel olduğu basitleştirilmiş durumu düşünün. $\mSat{M}{\Box!B}[w]$ ancak
  ve ancak her $v\in W$ için $\mSat{M}{!B}[v]$; $\mModel{M^*}$ için de ancak
  ve ancak $W^*$ içindeki her $[v]$ için $\mSat{M^*}{!B}[{[v]}]$ geçerli olsun. İlk
  fikir olarak, $u$ ile $v$ dünyaları $\mModel{M}$ içindeki
bütün !!{propositional variable}s konusunda uyuşuyorsa, yani
$\mSat{M}{p}[u]$ ancak ve ancak $\mSat{M}{p}[v]$ ise denk (aynı bölümde)
sayılsın. $V^*(p)=\Setabs{[w]}{\mSat{M}{p}[w]}$ olsun. Amacımız
$\mSat{M}{!A}[w]$ ancak ve ancak $\mSat{M^*}{!A}[{[w]}]$ göstermektir.
Bunu açıkça tümevarımla ispatlarız. Taban durumu $!A\equiv p$ olsun. Önce
  $\mSat{M}{p}[w]$ varsayalım. Tanım gereğince $[w]\in V^*(p)$; dolayısıyla
$\mSat{M^*}{p}[{[w]}]$. Tersine $\mSat{M^*}{p}[{[w]}]$ varsayalım. Bu,
$[w]\in V^*(p)$, yani $w$'ye denk bazı $v$ için $\mSat{M}{p}[v]$ demektir.
Fakat ``$w$, $v$'ye denk'' demek, ``$w$ ile $v$ tam olarak aynı
!!{propositional variable}s{}ni doğru yapar'' demektir; dolayısıyla
$\mSat{M}{p}[w]$. Tümevarım adımında örneğin $!A\ident\lnot!B$ olsun.
O zaman $\mSat{M}{\lnot!B}[w]$ ancak ve ancak $\mSat/{M}{!B}[w]$ ancak ve
ancak (tümevarım varsayımıyla) $\mSat/{M^*}{!B}[{[w]}]$ ancak ve ancak
$\mSat{M^*}{\lnot!B}[{[w]}]$. Modal olmayan diğer işleçler için de
benzerdir. $\Box$ için de işler: $\mSat{M^*}{\Box!B}[{[w]}]$ varsayalım.
Bu, her $[u]$ için $\mSat{M^*}{!B}[{[u]}]$ demektir. Tümevarım
varsayımıyla her $u$ için $\mSat{M}{!B}[u]$; dolayısıyla
$\mSat{M}{\Box!B}[w]$.

$\mModel{M^*}$ için erişilebilirlik bağıntısını da tanımlamamız gereken
genel durumda işler daha karmaşıktır. Yukarıdaki tümevarımlı ispatın
işlemesi için gereken koşulları erişilebilirlik bağıntısı $R^*$ sağlıyorsa
bir $\mModel{M^*}$ modeline \emph{süzme} diyeceğiz. O zaman her
$\mModel{M^*}$ süzmesi, $!A$'yı $[w]$'de ancak ve ancak $\mModel{M}$
$!A$'yı $w$'de doğru yapıyorsa doğru yapar. Fakat artık süzmelerin
\emph{var} olduğunu, yani $R^*$'yi gereken koşulları sağlayacak biçimde
  tanımlayabildiğimizi de göstermeliyiz. Bunun işlemesi için $u$ ve $v$ dünyalarını
yalnızca bütün !!{propositional variable}s konusunda değil, $!A$'nın bütün
alt !!{formula}s{}i konusunda uyuştuklarında denk saymalıyız. $!A$ yalnızca
sonlu sayıda alt !!{formula}s içerdiğinden bu yine süzmenin sonlu olmasını
güvence altına alır. Bir süzmeyi tanımlamanın tek bir yolu yoktur ve
süzmenin erişilebilirlik bağıntısının gerekli özellikleri (yansımalılık,
geçişlilik vb.) sağlaması için $R^*$'nin tanımında yaratıcı olmamız gerekir.

